% ============================================================================
% LANGENTWURF LE-003B: ``Wo ist?'' & Geräte-Sicherheit
% Profil: S (Senioren 65+) | Tier: 2 (Standard)
% ============================================================================

\documentclass[11pt,a4paper]{article}
\usepackage[utf8]{inputenc}
\usepackage[T1]{fontenc}
\usepackage[ngerman]{babel}
\usepackage{geometry}
\usepackage{fancyhdr}
\usepackage{lastpage}
\usepackage{xcolor}
\usepackage{tcolorbox}
\usepackage{enumitem}
\usepackage{longtable}
\usepackage{booktabs}
\usepackage{hyperref}
\usepackage{amssymb}

\geometry{left=2.5cm, right=2.5cm, top=2.5cm, bottom=2.5cm}
\definecolor{fach}{HTML}{b35c00}
\definecolor{gray}{HTML}{666666}
\definecolor{lightgray}{HTML}{f5f5f5}
\definecolor{successgreen}{HTML}{2a7a4b}
\definecolor{findmygreen}{HTML}{34C759}

\tcbuselibrary{skins,breakable}
\newtcolorbox{infobox}[1][]{ colback=lightgray, colframe=fach, fonttitle=\bfseries, title=#1, breakable }
\newtcolorbox{scorebox}{ colback=white, colframe=fach, fonttitle=\bfseries, title=Didaktik-Score, breakable }

\pagestyle{fancy}
\fancyhf{}
\fancyhead[L]{\small\textcolor{gray}{Digitale Grundlagen -- 003B ``Wo ist?''}}
\fancyhead[R]{\small\textcolor{gray}{LE Tier 2 / Profil S}}
\fancyfoot[C]{\small\thepage{} / \pageref{LastPage}}
\renewcommand{\headrulewidth}{0.4pt}

\begin{document}

\begin{titlepage}
    \centering
    \vspace*{2cm}
    {\Large\textcolor{gray}{Digitale Grundlagen -- Senioren 65+}}\\[2cm]
    {\Huge\textcolor{fach}{\textbf{Langentwurf}}}\\[0.5cm]
    {\large Tier 2 | Profil S}\\[2cm]
    {\LARGE\textbf{003B -- ``Wo ist?'' \& Geräte-Sicherheit}}\\[0.5cm]
    {\large Modul A: Geräte \& Zugänge}\\[2cm]
    \begin{tabular}{rl}
        \textbf{Zeitrahmen:} & 45--60 Minuten \\
        \textbf{Vorkenntnisse:} & Einheit 003 (Apple-ID, iCloud) \\
    \end{tabular}
    \vfill
    {\normalsize Stand: \today}
\end{titlepage}

\tableofcontents
\newpage

\section{Bedingungsanalyse}

\subsection{Ausgangssituation}
\begin{itemize}
    \item Teilnehmer kennen die Apple-ID und iCloud aus Einheit 003
    \item ``Wo ist?'' wurde in 003 nur erwähnt, nicht erklärt
    \item Viele wissen nicht, dass ``Wo ist?'' vorher aktiviert werden muss
    \item Große Angst vor Geräteverlust, aber keine Vorbereitung
\end{itemize}

\subsection{Typische Situation}
\begin{itemize}
    \item ``Ich habe mein iPhone verloren -- was kann ich tun?''
    \item ``Kann mir jemand helfen, es zu finden?''
    \item ``Ich wusste gar nicht, dass man das aktivieren muss!''
    \item ``Wie funktioniert das mit dem Computer?''
\end{itemize}

\section{Sachanalyse}

\subsection{Kernkonzepte}
\begin{enumerate}
    \item \textbf{``Wo ist?''-Funktion:} Alle Apple-Geräte auf einer Karte lokalisieren
    \item \textbf{Aktivierung erforderlich:} Muss VOR dem Verlust eingerichtet sein
    \item \textbf{Drei wichtige Schalter:} Mein iPhone suchen, Netzwerk ``Wo ist?'', Letzten Standort senden
    \item \textbf{Offline-Ortung:} Funktioniert über Bluetooth-Netzwerk anderer Apple-Geräte
    \item \textbf{Handlungsoptionen:} Ton abspielen, Als verloren markieren, Gerät löschen
    \item \textbf{Browser-Zugang:} icloud.com/find für Zugriff ohne Apple-Gerät
\end{enumerate}

\subsection{Alltagsanalogien}
\begin{tabular}{|l|p{8cm}|}
\hline
\textbf{Konzept} & \textbf{Analogie} \\
\hline
``Wo ist?'' & GPS-Tracker wie für Schlüssel (Tile, AirTag) \\
Aktivierung & Alarmanlage: Muss eingeschaltet sein, bevor eingebrochen wird \\
Ton abspielen & Schlüsselfinder: Piepsen lassen \\
Als verloren markieren & Sperren und Nachricht hinterlassen \\
Gerät löschen & Selbstzerstörung -- alle Daten weg \\
icloud.com/find & Fernbedienung über den Computer \\
\hline
\end{tabular}

\section{Didaktische Analyse}

\subsection{Stellung in der Reihe}
Einheit 003B vertieft das in 003 nur erwähnte ``Wo ist?''. Praktische Sicherheitsfunktion, die viele Senioren nicht kennen.

\subsection{Legitimation}
\begin{itemize}
    \item Geräteverlust ist reale Angst der Zielgruppe
    \item Präventive Einrichtung verhindert späteren Datenverlust
    \item Selbstständigkeit: Hilfe holen können, ohne Panik
    \item Vertrauenspersonen einbeziehen (wichtig für Senioren)
\end{itemize}

\subsection{Didaktische Reduktion}
\textbf{Ausgeklammert:} AirTags im Detail, Familienfreigabe-Ortung, Datenschutzdetails.

\textbf{Fokus:} ``Wo ist?'' aktivieren, verstehen, im Notfall nutzen.

\section{Lernziele}

\begin{tabular}{|c|p{7cm}|c|c|}
\hline
\textbf{Nr.} & \textbf{Die Teilnehmer können...} & \textbf{Stufe} & \textbf{Niveau} \\
\hline
FZ1 & erklären, was ``Wo ist?'' macht & Verstehen & Basis \\
FZ2 & die drei wichtigen Schalter in den Einstellungen finden & Anwenden & Basis \\
FZ3 & prüfen, ob ``Wo ist?'' auf ihrem Gerät aktiviert ist & Anwenden & Basis \\
FZ4 & ihr Gerät in der ``Wo ist?''-App auf der Karte finden & Anwenden & Basis \\
FZ5 & erklären, was ``Als verloren markieren'' bedeutet & Verstehen & Basis \\
FZ6 & icloud.com/find am Computer nutzen & Anwenden & Vertiefung \\
\hline
\end{tabular}

\section{Methodische Analyse}

\subsection{Vorgehen}
\begin{enumerate}
    \item \textbf{Einstieg:} ``Was machst du, wenn dein iPhone weg ist?'' -- Sammeln von Reaktionen
    \item \textbf{Problemstellung:} ``Aber ``Wo ist?'' funktioniert nur, wenn...''
    \item \textbf{Aktivierung:} Gemeinsam die drei Schalter finden und aktivieren
    \item \textbf{App ausprobieren:} Eigenes Gerät auf Karte finden, Ton abspielen
    \item \textbf{Notfall-Szenario:} Was tun bei echtem Verlust?
    \item \textbf{Notfall-Karte:} Zum Mitnehmen und Aufbewahren
\end{enumerate}

\subsection{Material}
Das vorhandene Material enthält:
\begin{itemize}
    \item Seite 1: Was ist ``Wo ist?'', Aktivierung, Verlust-Schritte
    \item Seite 2: Übungen (R/F, Situationen, praktische Übungen am Gerät)
    \item Seite 3: Checkliste + ausschneidbare Notfall-Karte
\end{itemize}

\section{Verlaufsplanung}

\begin{longtable}{|p{1cm}|p{2cm}|p{5cm}|p{3cm}|p{2cm}|}
\hline
\textbf{Zeit} & \textbf{Phase} & \textbf{Inhalt} & \textbf{TN-Aktivität} & \textbf{Material} \\
\hline
0--5 & Einstieg & ``Was machst du, wenn dein iPhone weg ist?'' & Nachdenken, Austausch & -- \\
\hline
5--10 & Problem & ``Wo ist?'' funktioniert nur wenn vorher aktiviert! & Zuhören & S.1 \\
\hline
10--25 & Aktivierung & Drei Schalter gemeinsam finden und prüfen & Am Gerät mitmachen & S.1 \\
\hline
25--35 & App testen & ``Wo ist?''-App öffnen, eigenes Gerät finden, Ton abspielen & Ausprobieren & S.2 \\
\hline
35--45 & Notfall & Was tun bei echtem Verlust? icloud.com/find erklären & Zuhören, Fragen & S.1 \\
\hline
45--55 & Sicherung & Checkliste, Notfall-Karte ausfüllen & Abhaken, Ausfüllen & S.3 \\
\hline
\end{longtable}

\section{Lösungen}

\textbf{Richtig/Falsch:}
\begin{enumerate}
    \item F -- Muss vorher aktiviert sein
    \item R -- Ton abspielen funktioniert
    \item F -- Mit ``Netzwerk Wo ist?'' auch offline
    \item R -- icloud.com/find im Browser
    \item R -- Löschen = Daten weg
    \item R -- Aktivierungssperre schützt
\end{enumerate}

\textbf{Was machst du?}
\begin{enumerate}
    \item In der Wohnung verlegt: Ton abspielen
    \item Gestohlen: Löschen + Polizei informieren
    \item Im Zug liegen lassen: Als verloren markieren
    \item Familie zeigen wo ich bin: Standort teilen (andere Funktion)
\end{enumerate}

\section{Didaktik-Score}

\begin{scorebox}
\begin{tabular}{|c|l|c|c|p{3.5cm}|}
\hline
\# & Dimension & Gew. & Score & Notiz \\
\hline
1 & Differenzierung & 15\% & 9/10 & Browser-Übung als Vertiefung \\
2 & Taxonomie & 10\% & 9/10 & Verstehen + Anwenden dominant \\
3 & Alltagsrelevanz & 10\% & 10/10 & Geräteverlust = echte Angst \\
4 & Aktivierung & 10\% & 10/10 & Praktisches Aktivieren am Gerät \\
5 & Scaffolding & 10\% & 9/10 & Schritt-für-Schritt-Anleitung \\
6 & Feedback & 10\% & 9/10 & Checkliste + Notfall-Karte \\
7 & Sprachliche Klarheit & 10\% & 9/10 & Einfache Sprache, gute Analogien \\
8 & Motivation & 10\% & 10/10 & ``Das brauchst du wirklich!'' \\
9 & Barrierefreiheit & 10\% & 9/10 & Große Schrift, klare Struktur \\
10 & Praktikabilität & 5\% & 9/10 & 55 Min gut machbar \\
\hline
\multicolumn{3}{|r|}{\textbf{GESAMT:}} & \textbf{9.30/10} & \textcolor{successgreen}{\textbf{FREIGABE}} \\
\hline
\end{tabular}
\end{scorebox}

\end{document}
